\documentclass[11pt]{article}
\usepackage[margin=1in]{geometry}
\usepackage{amsmath,amsfonts,amssymb,amsthm,enumerate,bbm}
\usepackage[]{graphicx}
\usepackage{color,subfigure}
\definecolor{scarlet}{rgb}{0.81,0,0}
\usepackage{multicol}
\usepackage{float}
\usepackage[all]{xypic}
\usepackage[colorlinks=true,citecolor=scarlet,linkcolor=scarlet]{hyperref}
\usepackage{colonequals}

\usepackage{fancyhdr, lastpage}
\pagestyle{fancy}
\fancyfoot[C]{{\thepage} of \pageref{LastPage}}



\DeclareMathOperator{\mSpec}{mSpec}
\DeclareMathOperator{\Spec}{Spec}
\DeclareMathOperator{\Ass}{Ass}
\DeclareMathOperator{\Supp}{Supp}
\DeclareMathOperator{\height}{height}
\DeclareMathOperator{\Hom}{Hom}
\DeclareMathOperator{\ann}{ann}
\DeclareMathOperator{\End}{End}
\DeclareMathOperator{\coker}{coker}
%\DeclareMathOperator{\ker}{ker}
\DeclareMathOperator{\rank}{rank}
\DeclareMathOperator{\im}{im}
\DeclareMathOperator{\M}{M}
\DeclareMathOperator{\Tor}{Tor}
\DeclareMathOperator{\id}{id}
\DeclareMathOperator{\ch}{char}
\DeclareMathOperator{\Aut}{Aut}
\DeclareMathOperator{\Gal}{Gal}
\DeclareMathOperator{\GL}{GL}

%\DeclareMathOperator{\dim}{dim}

\DeclareMathOperator{\lcm}{lcm}

\def\ra{\rightarrow}
\newcommand{\m}{\mathfrak{m}}
\newcommand{\C}{\mathbb{C}}
\newcommand{\Q}{\mathbb{Q}}
\newcommand{\Z}{\mathbb{Z}}
\newcommand{\ZZ}{\mathbb{Z}}
\newcommand{\R}{\mathbb{R}}
\newcommand{\N}{\mathbb{N}}
\newcommand{\ov}[1]{\overline{#1}}
\newcommand{\norm}{\trianglelefteq}

\def\ov#1{\overline{#1}}


\title{}
\date{\vspace{-0.5in}}

\makeatletter
\g@addto@macro\@floatboxreset\centering
\makeatother

\theoremstyle{definition}
\newtheorem{problem}{Problem}


\begin{document}

\thispagestyle{fancy}
\pagestyle{fancy}
\rhead{UNL $\mid$ Spring 2026}
\lhead{Introduction to Modern Algebra II}

\vspace{3em}

\begin{center}
	{\LARGE Some old qualifying exam questions\\}
\end{center}

\



\begin{problem} Consider $f(x) = x^6 + 3 \in \Q[x]$.

    \begin{enumerate}[(a)]
    \item Let\footnote{Hint: First show $\frac{a^3+1}2$ 
is primitive 6-th root of unity.} $a$ be a root of $f(x)$ and prove $\Q(a)$ is a Galois field extension of $\Q$. 

\item Find the Galois group $\Gal(\Q(a)/\Q)$.

     \end{enumerate}
     \end{problem}
     
     
     \begin{problem} Let $A$ be the matrix with entries in $\C$
     $$
     A  = \begin{bmatrix}
       3 & 0 & 0&  0 \\
       1 & 3 &0 &0\\
       1 &0 &3 &0 \\
       0&2&1&3\\
     \end{bmatrix}
         $$
         \begin{enumerate}[(a)]
         \item Find the Jordan canonical form of $A$.
         \item Is $A$ similar to
           $$
           \begin{bmatrix}
           0& -9& 0& 0 \\
           1& 6& 0& 0 \\
           0& 0& 0& -9 \\
           0&0&1&6 \\
         \end{bmatrix}?
         $$

     
       \end{enumerate}
       \end{problem}
       
       
       \begin{problem} Suppose $F \subseteq K$ is a finite extension of fields such that $[K: F]$ is odd. Show that for any $b \in K$ we have $F(b) = F(b^2)$.
\end{problem}


\begin{problem} Find, with justification, a complete and non-redundant list of conjugacy class representatives for the group $\GL_3(\mathbb{F}_2)$, where $\mathbb{F}_2$ is the field with two elements.
\end{problem}

\begin{problem}
  Let $p$ be any positive prime integer.
  \begin{enumerate}[(a)]
  \item Prove that if $p = k^2 + 1$ for some integer $k$, then $p$ is not an irreducible element of $\Z[i]$.
  
\item Prove $x^4 - p$ is irreducible in $\Q(i)[x]$. {\em Tip:} Find the degree of the splitting field of $x^4 - p$ over $\Q(i)$.  
\end{enumerate}
\end{problem}

\begin{problem} Let $E/F$ be a field extension and $\sigma:E\to E$ a field homomorphism fixing $F$; i.e., $\sigma(f)=f$ for all $f\in F$.
\begin{enumerate}[(a)]
\item Prove that if $E$ is algebraic over $F$ then $\sigma$ is an isomorphism.
\item Give an example of a specific field extension $E/F$ and a homomorphism $\sigma:E\to E$ fixing $F$ that is not an isomorphism.
\end{enumerate}
\end{problem}

\begin{problem} Let $L/F$ be a finite Galois extension with Galois group $G$. Show that $L=F(\alpha)$ if and only if the images of $\alpha$ under $G$ are distinct.
\end{problem}


\end{document}